\chapter{Manual de Referencia Técnica MDA-Audio-CIM-Machine}

Este documento es la especificación técnica completa del lenguaje \textbf{MDA-Audio-CIM-Machine}. Define la estructura de las máquinas de audio en el nivel más abstracto, así como las relaciones de sus elementos internos entre sí, y las entradas y salidas desde el exterior que pueden recibir o enviar. Siguiendo la nomenclatura de \textit{Model-Driven-Architecture} (MDA), este nivel se conoce como \textbf{CIM} (\textit{Computational Independence Model}).

En el CIM no se tiene porqué hablar de implementaciones concretas (osciladores, LFOs, filtros específicos) a bajo nivel, sino de \textbf{elementos} abstractos que capturan la intención del diseño de síntesis de la máquina.

\section{Bases del Modelo}

Una máquina es un conjunto de elementos que trabajan juntos para producir sonido. A su vez, como veremos más adelante, las máquinas cooperan entre sí para intercambiar sonido o valores de modificación (ejemplos: volumen, frecuencia, etc.) entre ellas.

\subsection{Modelo Raíz (\texttt{Document})}

El objeto raíz de un archivo MDA-Audio-CIM identifica a la máquina y contiene sus elementos compositivos.

\begin{table}[H]
\centering
\small
\begin{tabular}{@{}llp{4cm}l@{}}
\toprule
Atributo & Tipo & Descripción & Restricciones \\ \midrule
\texttt{id} & STRING & Identificador único de la máquina CIM. & Obligatorio. 36 car. (UUID). \\
\texttt{name} & STRING & Nombre descriptivo de la máquina. & Obligatorio. 1 a 20 car. \\
\texttt{description} & STRING & Propósito y contexto de la máquina. & Obligatorio. 10 a 600 car. \\
\texttt{elements} & ARRAY[Element] & Componentes que representan generadores o procesadores. & Obligatorio. Puede estar vacío. \\ \bottomrule
\end{tabular}
\end{table}

\textbf{Ejemplo de Documento:}
\begin{lstlisting}[language=json]
{
  "id": "550e8400-e29b-41d4-a716-446655440000",
  "name": "Oscillator Core",
  "description": "A fundamental oscillator machine for basic wave generation and shaping.",
  "elements": [ ... ]
}
\end{lstlisting}

\subsection{Interfaz del Elemento (\texttt{Element})}

Todos los bloques funcionales se definen como \texttt{Element}. Todos los campos base son obligatorios.

\begin{table}[H]
\centering
\small
\begin{tabular}{@{}llp{3cm}l@{}}
\toprule
Atributo & Tipo & Descripción & Restricciones \\ \midrule
\texttt{id} & STRING & Identificador único del componente. & Obligatorio. 36 car. (UUID). \\
\texttt{name} & STRING & Nombre del elemento. & Obligatorio. 1 a 20 car. \\
\texttt{description} & STRING & Descripción funcional. & Obligatorio. 10 a 600 car. \\
\texttt{params} & STRING & Parámetros modificables. & Obligatorio. 10 a 600 car. \\
\texttt{externalOutput} & ExternalOutput & Salida externa. & Obligatorio. \\
\texttt{externalInput} & ExternalInput & Entrada externa. & Obligatorio. \\
\texttt{sendTo} & ARRAY[SendTo] & Conexiones internas. & Obligatorio. Puede estar vacío. \\ \bottomrule
\end{tabular}
\end{table}

\textbf{Ejemplo de Elemento:}
\begin{lstlisting}[language=json]
{
  "id": "771e8400-e29b-41d4-a716-446655440111",
  "name": "Sine Generator",
  "description": "Generates a pure sine wave at a fixed base frequency.",
  "params": "frequency: 440Hz, gain: 0.8, phase: 0.0",
  "externalInput": { "hasExternalInput": false },
  "externalOutput": {
    "hasExternalOutput": true,
    "description": "Audio output of the sine wave."
  },
  "sendTo": []
}
\end{lstlisting}

\subsubsection{Configuración Externa}

Los subobjetos \texttt{externalOutput} y \texttt{externalInput} definen si el elemento sirve como uno de los puertos de entrada/salida global de la máquina.

\textbf{ExternalOutput / ExternalInput:}
\begin{itemize}
    \item \texttt{hasExternalOutput} / \texttt{hasExternalInput} (BOOLEAN): Habilita el puerto.
    \item \texttt{description} (STRING): Describe la naturaleza del puerto. Máximo 600 caracteres.
\end{itemize}

\subsection{Objeto de Conexión Interna (\texttt{SendTo})}

Define un vínculo direccional desde un elemento hacia otro dentro de la misma máquina. El elemento referenciado en el campo \texttt{idRef} debe existir dentro del documento.

\begin{table}[H]
\centering
\small
\begin{tabular}{@{}llp{4cm}l@{}}
\toprule
Atributo & Tipo & Descripción & Restricciones \\ \midrule
\texttt{id} & STRING & Identificador único de esta conexión. & Obligatorio. 36 car. (UUID). \\
\texttt{idRef} & REF[Element] & Referencia al elemento de destino. & Obligatorio. Debe existir. \\
\texttt{description} & STRING & Justificación técnica del vínculo. & Obligatorio. 10 a 600 car. \\ \bottomrule
\end{tabular}
\end{table}

\textbf{Ejemplo de Conexión:}
\begin{lstlisting}[language=json]
{
  "id": "b22e8400-e29b-41d4-a716-446655440333",
  "idRef": "c33e8400-e29b-41d4-a716-446655440444",
  "description": "Route filtered signal to the high pass filter element."
}
\end{lstlisting}

\begin{tcolorbox}[colback=red!5!white,colframe=red!75!black,title=IMPORTANTE]
Un elemento \textbf{no puede referenciarse a sí mismo} en su lista \texttt{sendTo}. El compilador rechaza cualquier auto-referencia como error fatal.
\end{tcolorbox}

\section{Tipos de Elemento de Audio (\texttt{Element})}

Representan cualquier bloque funcional de señal sonora dentro de la máquina. A este nivel no se especifica el algoritmo de síntesis: eso corresponde al PIM.

Los elementos típicos abarcan:
\begin{itemize}
    \item \textbf{Fuentes / Generadores:} Osciladores, samplers, generadores de ruido.
    \item \textbf{Modificadores / Efectos:} Filtros, reverbs, delays, envolventes ADSR, etc.
\end{itemize}

El campo \texttt{params} describe en lenguaje natural los parámetros conceptuales del elemento (ej.: \textit{``Frequency, Waveform''} o \textit{``Cutoff, Resonance''}).

\section{Reglas de Validación}

Se aplican las siguientes restricciones para garantizar que el modelo sea procesable:

\subsection{Reglas de Unicidad}

\begin{table}[H]
\centering
\small
\begin{tabular}{@{}lp{10cm}@{}}
\toprule
Alcance & Regla \\ \midrule
\textbf{Global} & Todos los \texttt{id} de todos los componentes deben ser únicos en el documento entero. \\
\textbf{Local} & No puede haber dos entradas en el array \texttt{sendTo} de un mismo componente con el mismo \texttt{id}. \\ \bottomrule
\end{tabular}
\end{table}

\subsection{Restricciones de Campos}

\begin{table}[H]
\centering
\small
\begin{tabular}{@{}lll@{}}
\toprule
Entidad & Campo & Restricción \\ \midrule
\texttt{Document} & \texttt{id} & Exactamente 36 caracteres. \\
\texttt{Document} & \texttt{name} & Entre 1 y 20 caracteres. \\
\texttt{Document} & \texttt{description} & Entre 10 y 600 caracteres. \\
\texttt{Element} & \texttt{id} & Exactamente 36 caracteres. \\
\texttt{Element} & \texttt{name} & Entre 1 y 20 caracteres. \\
\texttt{Element} & \texttt{description} & Entre 10 y 600 caracteres. \\
\texttt{Element} & \texttt{params} & Entre 10 y 600 caracteres. \\
\texttt{Element} & \texttt{externalDesc} & Máximo 600 caracteres. \\
\texttt{SendTo} & \texttt{id} & Exactamente 36 caracteres. \\
\texttt{SendTo} & \texttt{description} & Entre 10 y 600 caracteres. \\ \bottomrule
\end{tabular}
\end{table}

\subsection{Restricciones de Conectividad}

\begin{table}[H]
\centering
\small
\begin{tabular}{@{}lp{10cm}@{}}
\toprule
Regla & Descripción \\ \midrule
\textbf{Sin auto-referencia} & El \texttt{idRef} de un \texttt{SendTo} no puede apuntar al mismo elemento que lo contiene. \\
\textbf{Referencia válida} & El \texttt{idRef} debe resolver a un \texttt{Element} existente en el documento. \\ \bottomrule
\end{tabular}
\end{table}

\section{Guía de Uso}

CIM define conceptualmente una máquina de audio, entendiendo esta como un módulo que puede trabajar en solitario o bien en conjunto con otras máquinas (véase el documento de relaciones entre máquinas CIM). Proporciona el marco para que los diseñadores de sonido establezcan las capacidades de entrada/salida y el flujo interno sin preocuparse por la implementación técnica subyacente. Los modelos resultantes son la base para generar las especificaciones PIM detalladas.
